\documentclass[11pt, twocolumn, a4paper]{article}
\usepackage[text={186mm, 260mm},centering]{geometry}
\usepackage[utf8]{inputenc}
\usepackage[czech, shorthands=off]{babel}
\usepackage[T1]{fontenc}
\usepackage{lmodern}
\usepackage{amsmath}
\usepackage{amsthm}
\usepackage{amssymb}
\usepackage[hidelinks]{hyperref}
\hypersetup{hidelinks}


\theoremstyle{plain}
\newtheorem{defin}{Definice}

\theoremstyle{plain}
\newtheorem{veta}{Věta}

\begin{document}
\begin{titlepage}
    \begin{center}
        \Huge
        \textsc{Vysoké učení technické v~Brně\\
        \huge{Fakulta informačních technologií}}\\ \vspace{\stretch{0.382}}     
        {\huge{Typografie a~publikování\,--\,2. projekt\\
        Sazba dokumentů a~matematických výrazů}}\\ \vspace{\stretch{0.618}} 
    \end{center}
    {\Large2025\hfill Matúš Fignár (xfignam00)}
\end{titlepage}

\section*{Úvod}
V~této úloze vysázíme titulní stranu a~ukázku matematického textu,
v~němž se vyskytují například rovnice~\eqref{eq_7} na straně~\pageref{pg_one}, Věta~\ref{lem_1} nebo Definice~\ref{def_2}.
Pro vytvoření těchto odkazů používáme kombinace příkazů
\verb|\label|, \verb|\ref|, \verb|\eqref| a~\verb|\pageref|.
Před odkazy patří nezlomitelná mezera.
Text zvýrazníme pomocí příkazu \verb|\emph|, strojopisné písmo pomocí \verb|\texttt|.
Pro \LaTeX{}ové příkazy (s obráceným lomítkem) použijeme \verb|\verb|.

Titulní strana je vysázena prostředím \texttt{titlepage} a~nadpis je v~optickém středu
s~využitím \emph{zlatého řezu}, který byl probrán na přednášce.
Na titulní straně jsou tři různé velikosti písma a~mezi dvojicemi řádků textu
je řádkování se zadanou  velikostí 0{,}5\,em a~0{,}6\,em\footnote{Použijte správnou velikost mezery mezi číslem a~jednotkou.}


\section{Matematický text}
Symboly číselných množin sázíme makrem \verb|\mathbb|,
kaligrafická písmena  makrem \verb|\mathcal|.
Pozor na tvar i~sklon řeckých písmen: srovnejte \verb|\rho| a~\verb|\varrho|.
Konstrukce \verb|${}$| nebo \verb|\mbox{}| zabrání zalomení výrazu.

Pro definice a~věty slouží prostředí definovaná příkazem \verb|\newtheorem| z~balíku \texttt{amsthm}.
Tato prostředí obracejí význam \verb|\emph|:
uvnitř textu sázeného kurzívou se zvýrazňuje písmem v~základním řezu.
Důkazy se někdy ukončují značkou \verb|\qed|.

\subsection{Pseudometrický prostor}
Pro zarovnání rovností a nerovnosti pod sebe použijte vhodné prostředí.

\begin{defin}
V~\emph{pseudometrickém prostoru} $\mathcal{M} = (M, \varrho)$ 
značí $M$ množinu bodů, $\varrho : M \times M \rightarrow \mathbb{R} $ 
je zobrazení zvané \emph{pseudometrika}, které pro každé 
body $ x, y, z \in M $ splňuje následující podmínky:
\setcounter{equation}{0}
\begin{eqnarray}
    \varrho(x, x) & = & 0 \label{eq_1} \\ 
    \varrho(x, y) & = & \varrho(y, x) \label{eq_2} \\ 
    \varrho(x, y) + \varrho(y, z) & \geq & \varrho(x, z) \label{eq_3} 
\end{eqnarray}
    
\end{defin}

\subsection{Metrika}
Funkční hodnota pseudometriky $\varrho$ se nazývá \emph{vzdálenost}.
Vzdálenost každých dvou bodů je nezáporná.

\begin{veta}
    \label{lem_1}
    Pro každé dva body $ x, y \in M$ pseudometrického prostoru $(M, \varrho)$ 
    platí $\varrho(x, y) \geq 0$.
\end{veta}

Důkaz: Nechť $x,\,y \in M$ a~označme $d = \varrho(x, y)$. Využitím~\eqref{eq_2} máme 
$ 2d = \varrho(x, y) + \varrho(y, x)$, z~nerovnosti~\eqref{eq_3} vyplývá $2d \geq \varrho(x,x)$
a z rovnosti~\eqref{eq_1} dostaneme $2d \geq \varrho(x, x) = 0$. Odtud plyne $d \geq 0$.\qed

Speciálním případem pseudometrických prostorů jsou prostory metrické,
v~nichž dva různé body mají vždy kladnou vzdálenost.

\begin{defin}
    \label{def_2}Nechť $\mathcal{M} = (M, \varrho)$ je pseudometrický prostor, 
    v~němž platí $\varrho(x, y) > 0$, kdykoliv $ x \neq y$. Potom $ \mathcal{M} $ se nazývá \emph{metrický prostor}
    a~$\varrho$ je jeho \emph{metrika}. 
\end{defin}

\section{Rovnice}
Velikost závorek a~svislých čar je potřeba přizpůsobit jejich obsahu.
K~tomu jsou určeny modifikátory \verb|\left| a~\verb|\right|.
\begin{eqnarray}
    \lim\limits_{p \rightarrow 0} \left( \frac{1}{n} \sum\limits_{i=1}^n x_i^p \right)^{\frac{1}{p}}&=&\left( \prod\limits_{i=1}^{n} x_i \right)^{\frac{1}{n}}
\end{eqnarray}

Zde vidíme, jak se vysází proměnná určující limitu v~běžném textu: $\lim_{m \rightarrow \infty} f(m)$.
Podobně je to i~s~dalšími symboly jako $\bigcup_{N\in\mathcal{M}}N$ či $\sum^{m}_{i=1} x_{i}^{2}$.
S~vynucením méně úsporné sazby příkazem \verb|\limits| budou vzorce vysázeny v~podobě $\lim\limits_{m \rightarrow \infty} f(m)$ a~$ \sum\limits^{m}_{i=1} x_{i}^{2}$.
Složitější matematické formule sázíme mimo plynulý text pomocí prostředí \texttt{displaymath}.
\begin{eqnarray}
    \lim\limits_{n\rightarrow\infty}\left(1 + \frac{x}{n}\right)^{n} & = & \sum\limits_{n=0}^{\infty}\frac{x^n}{n!}\\
    \sum\limits_{\emptyset\neq X \subseteq P}(-1)^{\left| X \right| - 1}\left| \bigcap X\right| & = & \left| \bigcup P \right|\\
    -\int_{a}^{b} f(x)\,\mathrm{d}x & = & \int_{b}^{a} f(y)\,\mathrm{d}y \label{eq_7}
\end{eqnarray}
Nezapomeňte rovnice, na které se odkazujete, označit vhodným jménem pomocí \verb|\label|.

\section{Matice}
Pro sázení matic se používá prostředí \texttt{array} a~závorky s~výškou nastavenou pomocí
\verb|\left|, \verb|\right|.
$$
D = 
\left| \begin{array}{cccc}
    a_{11} & a_{12} & \cdots & a_{1n}\\
    a_{21} & a_{22} & \cdots & a_{2n}\\
    \vdots & \vdots & \ddots & \vdots\\
    a_{m1} & a_{m2} & \cdots & a_{mn}
        \end{array}\right|
= 
\left| \begin{array}{cc}
    x & y\\
    t & w
        \end{array}\right|
= xw - yt$$

Prostředí \texttt{array} lze úspěšně využít i~jinde,
například na pravé straně následující definiční rovnosti.
$$ B_n = \left\{
\begin{array}{l l}
1 & \text{pro } n = 0 \\
\sum\limits^{n-1}_{k=0}\binom{n}{k} B_k & \text{pro } n \geq 1
\end{array} \right. $$
Jestliže sázíme jen levou složenou závorku, pak za párovým \verb|\right|
místo závorky píšeme tečku.
\label{pg_one}
\end{document}
